\documentclass[10.5pt,a4paper]{article}
\def\yo{A}
% Macros y tabla de orden compartidas por los tres guiones. \yo = A, B o C.
\usepackage[utf8]{inputenc}
\usepackage[T1]{fontenc}
\usepackage{lmodern}
\usepackage[spanish,es-nodecimaldot]{babel}
\usepackage[a4paper,margin=2cm]{geometry}
\usepackage{amsmath,amssymb}
\usepackage[table]{xcolor}
\usepackage{booktabs,longtable,array}
\usepackage{fancyhdr}
\usepackage{enumitem}
\usepackage[hidelinks]{hyperref}
\setlist{itemsep=1pt,topsep=2pt,parsep=0pt}
\setlength{\parindent}{0pt}
\setlength{\parskip}{4pt}
\definecolor{azulfila}{HTML}{DCE9F8}
\definecolor{hipo}{HTML}{FFF3D6}
\definecolor{nuevo}{HTML}{DFF3E3}
\newcommand{\dt}{\ensuremath{\Delta t}}
\newcommand{\dts}{\ensuremath{\Delta t^{*}}}
\newcommand{\vv}{\ensuremath{\langle v^4\rangle/\langle v^2\rangle^2}}
\newcommand{\kbt}{\ensuremath{k_BT}}
\newcommand{\tn}[2]{\ensuremath{#1\times10^{#2}}}
\newcommand{\pr}[1]{\par\textbf{P:} #1\par\nobreak}
\newcommand{\rp}[1]{\textbf{R:} #1\par}
\newcommand{\hipotesis}[1]{\par\colorbox{hipo}{\parbox{\dimexpr\linewidth-2\fboxsep}{\textbf{Hipótesis nuestra (no medida):} #1}}\par}
\newcommand{\nuevo}[1]{\par\colorbox{nuevo}{\parbox{\dimexpr\linewidth-2\fboxsep}{\textbf{Nuevo en la presentación:} #1}}\par}
\newcommand{\debil}[1]{\par\textbf{Punto débil:} #1\par}
\usepackage{xspace}
\newcommand{\pc}{\char37\xspace}

\pagestyle{fancy}
\fancyhf{}
\renewcommand{\headrulewidth}{0pt}
\newcommand{\mi}[1]{\cellcolor{azulfila}\textbf{#1}}
\newcommand{\Fmi}[5]{\mi{#1} & \mi{#2} & \mi{#3} & \mi{#4} & \mi{#5} \\}
\newcommand{\Fno}[5]{#1 & #2 & #3 & #4 & #5 \\}
\newcommand{\F}[5]{\ifnum\pdfstrcmp{#3}{\yo}=0 \expandafter\Fmi\else\expandafter\Fno\fi{#1}{#2}{#3}{#4}{#5}}
\newcommand{\interv}[8]{%
  \par\vspace{6pt}\noindent{\large\bfseries Diapositiva #1 \textemdash{} #2}\hfill{\small reloj #3 \textbullet{} #4 s}\par\nobreak
  {\small\textit{Entrás después de:} #5}\par\nobreak
  #6\par
  \ifx\relax#7\relax\else\colorbox{hipo}{\parbox{\dimexpr\linewidth-2\fboxsep}{\small\textbf{Números a mano:} #7}}\par\fi
  {\small\textit{#8}}\par}
\newcommand{\nota}[1]{\colorbox{nuevo}{\parbox{\dimexpr\linewidth-2\fboxsep}{\small #1}}\par}
\newcommand{\cabecera}[3]{%
  {\LARGE\bfseries Guion Orador #1}\par
  {\large TP4 \textemdash{} Dinámica molecular regida por el paso temporal \textemdash{} Presentación de 13 minutos}\par\vspace{2pt}
  \textbf{Tu tiempo hablando:} #2 en #3 intervenciones. \textbf{Tiempo total planeado de la presentación:} 12:30 (límite 13:00, quedan 30 s de margen para pausas y pases).

  Los textos están pensados a unas 140--150 palabras por minuto. Si te quedás corto o largo, respetá los segundos de la tabla: el reloj de cada diapositiva es acumulado. Los números de diapositiva son los del pie «n / 33» de la presentación vigente.}
\newcommand{\tablaorden}{%
\section{Orden completo (tus filas en azul)}
{\small
\begin{longtable}{@{}rp{8.6cm}crr@{}}
\toprule Diap. & Contenido & Voz & Seg. & Reloj \\ \midrule \endhead
\F{1}{Título}{A}{10}{0:00}
\F{3}{Sistema real}{A}{20}{0:10}
\F{4}{Modelo matemático}{B}{25}{0:30}
\F{6}{Integración: Verlet original}{C}{25}{0:55}
\F{7}{Arquitectura del motor}{A}{25}{1:20}
\F{9}{Sistema simulado}{B}{25}{1:45}
\F{10}{Observables (1/2)}{C}{30}{2:10}
\F{11}{Observables (2/2)}{A}{30}{2:40}
\F{12}{Corridas realizadas}{B}{20}{3:10}
\F{14}{Oscilador: ECM contra $\Delta t$}{C}{30}{3:30}
\F{15}{Energía total en función del tiempo}{A}{20}{4:00}
\F{16}{Error de energía contra $\Delta t$}{B}{40}{4:20}
\F{17}{Tiempo de ejecución en función de $N$}{C}{40}{5:00}
\F{18}{Tiempo de ejecución en función de la densidad (nueva)}{B}{20}{5:40}
\F{19}{Variación de $x_0$ (animaciones)}{A}{20}{6:00}
\F{20}{Fracción de usadas en función del tiempo}{B}{25}{6:20}
\F{21}{Tiempo al 90\,\pc contra $x_0$ (parte 1)}{C}{30}{6:45}
\F{21}{Tiempo al 90\,\pc contra $x_0$ (parte 2)}{A}{20}{7:15}
\F{22}{Billar sin obstáculos (animación)}{A}{15}{7:35}
\F{23}{Estacionario: $k_BT$ ajustado en el tiempo (nueva)}{B}{25}{7:50}
\F{24}{Distribución de rapideces: evolución}{C}{25}{8:15}
\F{25}{Búsqueda del mejor ajuste}{A}{25}{8:40}
\F{26}{Distribución estacionaria y ajuste}{B}{30}{9:05}
\F{27}{Tiempos de conversión contra la densidad}{C}{30}{9:35}
\F{28}{$\langle t_{90}\rangle$ y $\langle t_{100}\rangle$ contra la densidad (nueva)}{A}{20}{10:05}
\F{29}{Fracción de usadas a $t_f$ contra la densidad}{A}{20}{10:25}
\F{30}{Mapa de $t_{90}$ en $(x_0,N)$ --- parte 1}{B}{25}{10:45}
\F{30}{Mapa de $t_{90}$ en $(x_0,N)$ --- parte 2}{C}{20}{11:10}
\F{32}{Conclusiones (1 a 3)}{A}{20}{11:30}
\F{32}{Conclusiones (4 a 6)}{B}{20}{11:50}
\F{32}{Conclusiones (7 y 8) y cierre}{C}{20}{12:10}
\bottomrule
\end{longtable}}

\textbf{Reparto.} Las diapositivas de resultados (la parte más importante) las hablamos los tres, intercaladas: A, B y C se turnan. En las diapositivas de tiempo al 90\,\pc contra $x_0$ (C y A) y del mapa de calor (B y C) hablan dos voces seguidas en la misma diapositiva; el pase se hace sin cambiar de lámina. Las diapositivas 14 a 18 no se detienen en «por qué»: las razones del método (Verlet, regla de $\Delta t^*$) ya no están escritas, se dicen de viva voz.

\textbf{Cambios respecto del guion anterior.} Las tres diapositivas nuevas son la 18, la 23 y la 28 (marcadas «nueva»); la antigua «Relajación hacia Maxwell-Boltzmann» (cociente) ya no tiene diapositiva propia: ahora su idea se dice en la 11 (definición) y en la 23 ($k_BT(t)$). Desde la 18 la numeración se corrió.}
\newcommand{\antes}{%
\section{Antes de subir}
Leé tu guion en voz alta con cronómetro: cada intervención debe entrar en sus segundos. Los pases son frases cortas (``sigue B'', ``te paso a C'') sin cambiar de tema.

Para las preguntas, mirá la \textbf{Guía de estudio TP4}: cada uno defiende sus diapositivas, pero cualquiera puede ser interrogado sobre cualquier parte.}

\cfoot{Guion Orador A \textbullet{} TP4 Grupo 5 \textbullet{} \thepage}
\begin{document}
\cabecera{A}{4:05}{12}
\tablaorden

\section{Tu guion}

\interv{1}{Título}{0:00}{10}{---}{Arrancás vos. Buenos días. Somos el Grupo 5 y les vamos a presentar el TP4: dinámica molecular regida por el paso temporal, aplicada a un billar circular.}{}{Seguís vos en la diapositiva 3 (sin cambiar de lámina de voz).}

\interv{3}{Sistema real}{0:10}{20}{vos mismo (título)}{Nuestro sistema real es un gas granular: partículas que sólo interactúan cuando se tocan. Los choques son blandos y, como no hay disipación, la energía se conserva. Lo simulamos con dinámica molecular regida por el paso temporal: el estado avanza con un $\Delta t$ fijo, integrando las ecuaciones de movimiento, sin calcular eventos.}{}{Pasás a B (Modelo matemático).}

\interv{7}{Arquitectura del motor}{1:20}{25}{C (Integración: Verlet original)}{El motor es C++20, sin dependencias y de un solo hilo. En cada paso armamos una grilla de celdas con el Cell Index Method, calculamos las fuerzas entre pares vecinos, contra los obstáculos y contra la pared con la partícula imagen, damos el paso de Verlet, marcamos las conversiones de fresca a usada y escribimos el estado cada $n$ pasos. El costo por paso es lineal en $N$.}{$M=29$ celdas por lado ($\ge2r$); pared: $x_{img}=(R+r)\hat n$}{Pasás a B (Sistema simulado).}

\interv{11}{Observables (2/2)}{2:40}{30}{C (Observables 1/2)}{Para la distribución de rapideces usamos bines de 0,05 metros por segundo, normalizados como densidad de probabilidad. Para saber cuándo relaja usamos el cociente $\langle v^4\rangle$ sobre $\langle v^2\rangle$ al cuadrado: vale 1 si todas tienen la misma rapidez y 2 en Maxwell-Boltzmann 2D. Ajustamos $k_BT$ con el método de la Teórica 0: barremos el parámetro, calculamos el error cuadrático contra el modelo y nos quedamos con el mínimo. La densidad es $N$ sobre $\pi R$ cuadrado.}{barrido $k_BT$: $10^{-3}$ a $5\times10^{-2}$ J, paso $10^{-5}$ J}{Pasás a B (Corridas realizadas).}

\interv{15}{Energía total en función del tiempo}{4:00}{20}{C (Oscilador: ECM contra $\Delta t$)}{Ahora elegimos $\Delta t$ para el billar: $N$ igual 300, sin obstáculos. Con $\Delta t$ de $10^{-3}$ la energía se multiplica por más de diez. Con $\Delta t$ chicos queda plana. Para decidir, resumimos cada curva en un escalar: el desvío relativo medio de la energía.}{$E(0)=Nmv_0^2/2=3{,}75$ J}{Pasás a B (Error de energía contra $\Delta t$).}

\interv{19}{Variación de $x_0$ (animaciones)}{6:00}{20}{B (Tiempo de ejecución contra la densidad)}{Entramos al estudio central: dónde poner los obstáculos. Mostramos dos animaciones con $N$ igual 100: $x_0$ igual a $r$, con los obstáculos tocándose, y $x_0$ igual a 0,20 metros. Cada fotograma es el instante en que hay 50 usadas. Los links a las animaciones están en la diapositiva.}{50 usadas: 4,5 s ($x_0=0{,}20$ m) vs 7,8 s ($x_0=r$)}{Pasás a B (Fracción de usadas en función del tiempo).}

\interv{21}{Tiempo al 90\,\pc contra $x_0$ (parte 2)}{7:15}{20}{C (misma diapositiva)}{Los mínimos vecinos no se distinguen dentro del desvío, por eso hablamos de zona favorable y no de un $x_0$ óptimo. En los extremos el blanco efectivo es menor: con los obstáculos juntos se tapan entre sí, y contra la pared queda inaccesible más de la mitad de su contorno; es una interpretación nuestra. Con $N$ igual 20 todo es más rápido, porque hay menos choques entre partículas que las desvíen.}{$N=20$: mínimo $(11\pm3)$ s en 0,40 m, fila plana entre 0,20 y 0,40}{Seguís vos en la diapositiva 22 (cambiás de lámina).}

\interv{22}{Billar sin obstáculos (animación)}{7:35}{15}{vos mismo (Tiempo al 90\,\pc)}{Pasamos al inciso 2.3: billar sin obstáculos, $N$ igual 100, todas con rapidez $v_0$. El fotograma es de $t$ igual 2 segundos. Sólo los choques entre partículas pueden repartir la energía, porque la pared conserva el módulo de la velocidad.}{}{Pasás a B (Estacionario: $k_BT$ ajustado en el tiempo).}

\interv{25}{Búsqueda del mejor ajuste}{8:40}{25}{C (Distribución de rapideces: evolución)}{Para ajustar usamos el método de la Teórica 0, con un único parámetro, $k_BT$. Calculamos el error cuadrático entre la $f$ medida y la de Maxwell-Boltzmann para cada $k_BT$, y nos quedamos con el mínimo. La curva tiene un único mínimo bien definido: $1{,}24\times10^{-2}$ joules, contra $mv_0^2/2$, que es $1{,}25\times10^{-2}$.}{$k_BT=(1{,}24\pm0{,}02)\times10^{-2}$ J; $mv_0^2/2=1{,}25\times10^{-2}$ J ($-1\,\pc$); estrella = mínimo, trazos = $mv_0^2/2$}{Pasás a B (Distribución estacionaria y ajuste).}

\interv{28}{$\langle t_{90}\rangle$ y $\langle t_{100}\rangle$ contra la densidad}{10:05}{20}{C (Tiempos de conversión contra la densidad)}{Mismo estudio, ahora con $x_0$ igual 0,20 y $t_{max}$ de 100 segundos. El $t_{90}$ crece con la densidad, de 12 segundos con $N$ igual 20 a 45 con $N$ igual 400. El $t_{100}$ también crece, y con $N$ igual 400 ninguna realización lo alcanza. No hay una densidad óptima distinguible.}{$\langle t_{90}\rangle$: 12, 14, 15, 17, 20, 25, 31, 45 s ($N=20,\dots,400$); $\langle t_{100}\rangle$: 21, 28, 35, 42, 58, 63 s hasta $N=250$; $N=300$ parcial (9 de 10, cota 79 s)}{Seguís vos en la diapositiva 29.}

\nota{La diapositiva 28 es otra condición que la 27 (RSA, $x_0=0{,}20$ m, $t_{max}=100$ s, mapa 2.4b; la 27 es red, $x_0=r$, $t_f=30$ s). Si te preguntan, está explicado en la Guía, Parte 6.}

\interv{29}{Fracción de usadas a $t_f$ contra la densidad}{10:25}{20}{vos mismo (diapositiva 28)}{Como $t_{90}$ no siempre se alcanza en 30 segundos, miramos la fracción de usadas a $t_f$: disminuye al aumentar la densidad. Son las mismas corridas de $x_0$ igual a $r$ de la diapositiva 27. Con más partículas el movimiento es más difusivo y cada una tarda más en llegar a un obstáculo. A $N$ igual 650 casi nada se convirtió, y eso incluye que arrancamos de una red ordenada.}{$F_u(30\text{ s})$: 0,93 ($N=50$); 0,018 ($N=650$)}{Pasás a B (Mapa de $t_{90}$ en $(x_0,N)$, parte 1).}

\interv{32}{Conclusiones (1 a 3)}{11:30}{20}{C (Mapa, parte 2)}{Resumimos. Primero: en el oscilador, Beeman, Verlet original y Velocity Verlet dan casi el mismo error, de orden 2, y elegimos Verlet; Euler predictor-corrector es mucho peor. Segundo: el $\Delta t$ elegido conserva la energía por debajo del umbral. Tercero: el costo del paso temporal crece casi lineal con $N$ y supera al método por eventos a partir de cierto $N$.}{}{Pasás a B (Conclusiones 4 a 6).}

\antes
\end{document}
